\section*{第 \ref{ch:g11:cancer} 章——基因组损伤与癌症}

\begin{solution}{exo:g11:cancer:1}
\omterm{def:g11:cancer:cancer}{癌症}：\omterm{def:g10:cells-common-unit:cell}{细胞}不受正常调控地分裂并侵入
邻近组织。\omterm{def:g11:cancer:cancer}{肿瘤}：由这样一个克隆形成的团块。
转移灶：由\omterm{def:g10:cells-common-unit:cell}{细胞}经血液或淋巴迁移，
在别处形成的新\omterm{def:g11:cancer:cancer}{肿瘤}。
\end{solution}

\begin{solution}{exo:g11:cancer:2}
一种\omterm{def:g10:universal-dna:gene}{基因}，其\omterm{def:g11:gene-expression:protein}{蛋白质}传递生长信号，推动\omterm{def:g10:cells-common-unit:cell}{细胞}进入
\omterm{def:g11:cell-cycle-mitosis:cycle}{细胞周期}；若该\omterm{def:g11:gene-expression:protein}{蛋白质}发生\omterm{def:g11:mutations:mutation}{突变}而持续活化，则此\omterm{def:g10:universal-dna:gene}{基因}成为
\omterm{prop:g11:cancer:genes}{原癌基因}。一个卡住的“加速器”会驱动\omterm{def:g10:cells-common-unit:cell}{细胞}不断增殖，
不论另一个\omterm{def:g10:universal-dna:gene}{等位基因}如何。
\end{solution}

\begin{solution}{exo:g11:cancer:3}
\omterm{def:g10:universal-dna:information}{DNA} 受损时，它使\omterm{def:g11:cell-cycle-mitosis:cycle}{细胞周期}停止，并引发修复或
\omterm{def:g10:cells-common-unit:cell}{细胞}自毁。一个正常工作的\omterm{def:g10:universal-dna:gene}{等位基因}仍能产生足够的\omterm{def:g11:gene-expression:protein}{蛋白质}，
起制动作用；只有两个\omterm{def:g10:universal-dna:gene}{等位基因}都丧失，制动才会失去。
\end{solution}

\begin{solution}{exo:g11:cancer:4}
紫外线（皮肤癌）、烟草焦油（肺癌）、电离
辐射、石棉、酒精、乳头瘤病毒（宫颈癌）。
\end{solution}

\begin{solution}{exo:g11:cancer:5}
其所有\omterm{def:g10:cells-common-unit:cell}{细胞}都来自最初发生首次\omterm{def:g11:mutations:mutation}{突变}的同一个\omterm{def:g10:cells-common-unit:cell}{细胞}，
因此它们共享该\omterm{def:g11:mutations:mutation}{突变}及随后发生的\omterm{def:g11:mutations:mutation}{突变}，并
带有相同的标记，如相同的失活 X \omterm{def:g11:cell-cycle-mitosis:chromosome}{染色体}。
\end{solution}

\begin{solution}{exo:g11:cancer:6}
40 岁时约 200，80 岁时约 2700：年龄增加 2 倍，
发病数增至约 13 倍——多个事件相乘，而非单一的比例性原因。
\end{solution}

\begin{solution}{exo:g11:cancer:7}
约 10 与 17。每天吸 20 支烟者：$17 \times 1\% = 17\%$，
即约六分之一。
\end{solution}

\begin{solution}{exo:g11:cancer:8}
失去 p53 后，\omterm{def:g10:cells-common-unit:cell}{细胞}在复制前不再暂停以修复受损的 \omterm{def:g10:universal-dna:information}{DNA}，
损伤严重时也不会自行凋亡：损伤在每次分裂时
都被复制为\omterm{def:g11:mutations:mutation}{突变}，
\omterm{def:g10:universal-dna:gene}{基因}组混乱的\omterm{def:g10:cells-common-unit:cell}{细胞}得以存活。
\end{solution}

\begin{solution}{exo:g11:cancer:9}
她体内每个\omterm{def:g10:cells-common-unit:cell}{细胞}都已缺少一个\omterm{def:g10:universal-dna:gene}{等位基因}，一次体\omterm{def:g10:cells-common-unit:cell}{细胞}
\omterm{def:g11:mutations:mutation}{突变}即可使另一个\omterm{def:g10:universal-dna:gene}{等位基因}丧失，而无需两次独立
\omterm{def:g11:mutations:mutation}{突变}。由于出生时已迈出一步，其余步骤
更早达到，且涉及更多\omterm{def:g10:cells-common-unit:cell}{细胞}。她的姐妹从同一
亲代继承该\omterm{def:g10:universal-dna:gene}{等位基因}的概率为 $1/2$。
\end{solution}

\begin{solution}{exo:g11:cancer:10}
这些药物无差别地杀伤分裂中的\omterm{def:g10:cells-common-unit:cell}{细胞}：毛根\omterm{def:g10:cells-common-unit:cell}{细胞}、
肠道上皮与骨髓的\omterm{def:g10:cells-common-unit:cell}{细胞}分裂迅速，同样受到
\omterm{def:g11:cancer:cancer}{肿瘤}的伤害——因而出现脱发、恶心，以及白\omterm{def:g10:cells-common-unit:cell}{细胞}减少。
\end{solution}

\begin{solution}{exo:g11:cancer:11}
停药后，新\omterm{def:g11:mutations:mutation}{突变}的来源即告终止，且上皮\omterm{def:g10:cells-common-unit:cell}{细胞}
逐渐被更替，从而清除了许多携带早期步骤的克隆。
但有些携带多个步骤的克隆仍会存留，且已存在的\omterm{def:g11:mutations:mutation}{突变}
并不会被消除；其风险仍高于不吸烟者。
\end{solution}

\begin{solution}{exo:g11:cancer:12}
$10^{12} \times 10^{-7} = 10^5$ 个\omterm{def:g10:cells-common-unit:cell}{细胞}获得第一次\omterm{def:g11:mutations:mutation}{突变}。
每个\omterm{def:g10:cells-common-unit:cell}{细胞}都会形成一个克隆，该克隆可能分裂数千次，因此第二次
\omterm{def:g11:mutations:mutation}{突变}并非只在一个\omterm{def:g10:cells-common-unit:cell}{细胞}中寻找，而是在整个克隆的
分裂中寻找，第三次\omterm{def:g11:mutations:mutation}{突变}亦然：各步骤使靶点数目成倍增加，
因而真实概率比 $10^{-21}$ 高出多个数量级。
\end{solution}

\begin{solution}{exo:g11:cancer:13}
病毒蛋白在不影响\omterm{def:g10:universal-dna:gene}{基因}的情况下解除制动：被感染
\omterm{def:g10:cells-common-unit:cell}{细胞}的表现如同 p53 已丧失，并积累\omterm{def:g11:mutations:mutation}{突变}，
从而完成整个过程。疫苗使免疫系统
在病毒感染宫颈之前将其清除；一旦\omterm{def:g10:cells-common-unit:cell}{细胞}已被
感染并发生转化，疫苗便不再起作用——因此
应在青春期早期接种疫苗。
\end{solution}

\begin{solution}{exo:g11:cancer:14}
息肉是已经历一两个步骤的克隆；切除息肉
即去除了可能发生其余步骤的\omterm{def:g10:cells-common-unit:cell}{细胞}，因此
\omterm{def:g11:cancer:cancer}{癌症}便不会形成。息肉发展为\omterm{def:g11:cancer:cancer}{癌症}需十至
二十年，因此每十年一次的筛查可在大多数息肉
完成这一过程之前将其检出。
\end{solution}

\begin{solution}{exo:g11:cancer:15}
发病率随年龄上升，是因为各步骤需要数十年，而
这些步骤发生在一生之中，且主要受暴露因素影响——仅烟草
一项即占\omterm{def:g11:cancer:cancer}{癌症}死亡的三分之一，感染、
酒精、日晒与饮食则占其余大部分。无论何时消除某种暴露，
都能减少尚未出现的\omterm{def:g11:mutations:mutation}{突变}；筛查则在最后一步
之前清除克隆。老年是这一解释被呈现的时候，
而不是它被写成的时候。
\end{solution}

\begin{solution}{pb:g11:cancer:1}
\textbf{1.} 这些\omterm{def:g11:mutations:mutation}{突变}存在于\omterm{def:g11:cancer:cancer}{肿瘤}中，而不存在于患者的
其他\omterm{def:g10:cells-common-unit:cell}{细胞}中，因此它们发生于体\omterm{def:g10:cells-common-unit:cell}{细胞}。种系\omterm{def:g11:mutations:mutation}{突变}
则会存在于每个\omterm{def:g10:cells-common-unit:cell}{细胞}中，包括健康组织。

\textbf{2.} \omterm{def:g11:cancer:cancer}{肿瘤}中的大多数\omterm{def:g11:mutations:mutation}{突变}由
苯并芘——烟草烟雾中的\omterm{def:g11:mutations:mutagen}{诱变剂}留下了其特征。

\textbf{3.} 无义\omterm{def:g11:mutations:mutation}{突变}：提前出现的终止\omterm{def:g11:gene-expression:code}{密码子}使
p53 截短，因而无法发挥作用。

\textbf{4.} p53 是刹车：一个正常的\omterm{def:g10:universal-dna:gene}{等位基因}仍能起制动作用，因此两个
\omterm{def:g10:universal-dna:gene}{等位基因}都必须失活（一个发生\omterm{def:g11:mutations:mutation}{突变}，另一个缺失）。受体则是油门：
一种持续活化的\omterm{def:g11:gene-expression:protein}{蛋白质}便会驱动\omterm{def:g10:cells-common-unit:cell}{细胞}增殖，
无论另一个\omterm{def:g10:universal-dna:gene}{等位基因}如何。

\textbf{5.} 受体\omterm{def:g10:universal-dna:gene}{基因}是油门（\omterm{prop:g11:cancer:genes}{原癌基因}
转变为癌\omterm{def:g10:universal-dna:gene}{基因}）；p53 是刹车（\omterm{prop:g11:cancer:genes}{抑癌基因}）。

\textbf{6.} $20 \times 365 \times 35 \approx \num{255000}$ 支香烟。

\textbf{7.} 每个\omterm{def:g10:cells-common-unit:cell}{细胞}约有 $\num{255000}$ 次替换——
是所发现 5000 次的五十倍。大多数损伤被修复，\omterm{def:g11:mutations:mutation}{突变}过多的
\omterm{def:g10:cells-common-unit:cell}{细胞}会死亡；5000 即存活下来的那部分。

\textbf{8.} 这些\omterm{def:g11:mutations:mutation}{突变}发生较早，发生于起始\omterm{def:g10:cells-common-unit:cell}{细胞}或其最初的
后代中，因此每个\omterm{def:g11:cancer:cancer}{肿瘤}\omterm{def:g10:cells-common-unit:cell}{细胞}都继承了它们；其余\omterm{def:g11:mutations:mutation}{突变}则
较晚在亚克隆中产生。

\textbf{9.} $10^{10} \approx 2^{33}$：33 次倍增，每次
需 100 天，约 9 年。

\textbf{10.} 若\omterm{def:g11:cancer:cancer}{肿瘤}从其起始\omterm{def:g10:cells-common-unit:cell}{细胞}生长至
形成用了 9 年或更久，则最后一个关键\omterm{def:g11:mutations:mutation}{突变}
约出现在吸烟 35 年中的第 25 年左右，或更早。

\textbf{11.} $17 \times 1\% = 17\%$.

\textbf{12.} 17 例\omterm{def:g11:cancer:cancer}{癌症}中，1 例无论如何都会发生，因此
每 100 例中有 16 例可归因于吸烟。

\textbf{13.} 不再新增\omterm{def:g11:mutations:mutation}{突变}；上皮自我更新，
大量携带早期步骤的克隆随时间脱落并被替代，
其余克隆所具有的步骤比继续积累时要少，
即少于若继续积累本可获得的步骤。

\textbf{14.} 若干罕见\omterm{def:g11:mutations:mutation}{突变}须在同一谱系中累积，
且该谱系随后须生长到可检测的规模：每个阶段都需
数年，合计需数十年。

\textbf{15.} 每个\omterm{def:g10:cells-common-unit:cell}{细胞}每支香烟产生一个\omterm{def:g11:mutations:mutation}{突变}，发生于约 $10^9$
个上皮\omterm{def:g10:cells-common-unit:cell}{细胞}中，经过 $\num{255000}$ 支香烟，共约 $10^{14}$ 个\omterm{def:g11:mutations:mutation}{突变}，
分散于整个肺部；其中少数会击中 p53 或\omterm{prop:g11:cancer:genes}{原癌基因}，
发生于已带有另一步骤的\omterm{def:g10:cells-common-unit:cell}{细胞}中。正是试验次数之多，
使小概率事件成为必然。

\textbf{16.} 吸烟者：$\num{27000}$ 例死亡，在 $\num{11.25e6}$ 人中，
约每 100\,000 人 240 例；不吸烟者：3000 例，在 $\num{33.75e6}$ 人中，
约每 100\,000 人 9 例。

\textbf{17.} 约 27 倍——与图中重度吸烟者的 17
至 25 倍为同一数量级；总体数据混合了轻度和重度
吸烟者以及已戒烟者。

\textbf{18.} 约每年 3000 例，即不吸烟者所占的份额。并非
立即如此：现有吸烟者和既往吸烟者体内已存在的\omterm{def:g11:mutations:mutation}{突变}，
还将在此后二三十年内引发\omterm{def:g11:cancer:cancer}{癌症}。

\textbf{19.} 对每位吸烟者进行筛查，可挽救约 20\% 的
$\num{27000}$ 例死亡，即每年约 5000 例；戒烟
则最终约可预防 $\num{24000}$ 例，为前者的五倍。

\textbf{20.} G \omterm{def:g11:mutations:mutation}{突变}为 T 的特征证明该\omterm{def:g11:cancer:cancer}{肿瘤}的\omterm{def:g11:mutations:mutation}{突变}来自
烟草烟雾；该\omterm{def:g11:cancer:cancer}{癌症}由一个癌\omterm{def:g10:universal-dna:gene}{基因}被激活、
p53 两个\omterm{def:g10:universal-dna:gene}{等位基因}均丢失而形成；不吸烟本可预防它。
\end{solution}
